\section{Exercise 2: Advance}
\subsection{Task and method}
Implement \texttt{[yn, yorigin] = advance(xn, xorigin, k)} for
\begin{equation}
  y(n)=x(n+k),\qquad k>0,
\end{equation}
where $k$ is an integer number of samples. Advancing the signal moves
each sample to the left by $k$ time units. A sample originally at $n_i$
therefore appears at $n_i-k$. As with delay, the stored amplitudes and
their order remain unchanged:
\begin{align}
  \texttt{yn} &= \texttt{xn},\\
  i-\texttt{yorigin} &= i-\texttt{xorigin}-k,\\
  \texttt{yorigin} &= \texttt{xorigin}+k.
\end{align}
The origin parameter increases because it is a reference vector position,
not the direction of the time shift.

\subsection{Source code}
\lstinputlisting[caption={Scilab advance function and example call},label={lst:advance}]{../2.sci}

\subsection{Example and expected result}
Use the same input sequence as in Exercise 1 to compare the two operations:
\begin{lstlisting}[numbers=none]
[yn, yorigin] = advance([1, -2, 3, 6], 3, 1);
\end{lstlisting}
The calculated output is
\[
  \texttt{yn}=[1,-2,3,6],\qquad \texttt{yorigin}=4.
\]
\begin{table}[H]
  \centering
  \caption{Sample coordinates before and after a one-sample advance}
  \label{tab:advance}
  \begin{tabular}{lrrrr}
    \toprule
    Vector position $i$ & 1 & 2 & 3 & 4\\
    \midrule
    Sample value & 1 & $-2$ & 3 & 6\\
    Input time $n_i$ & $-2$ & $-1$ & 0 & 1\\
    Output time $n_i-1$ & $-3$ & $-2$ & $-1$ & 0\\
    \bottomrule
  \end{tabular}
\end{table}

\subsection{Execution evidence}
\resultfigure{images/exercise2.png}
  {Input and advanced output for Exercise 2}
  {fig:advance}
The captured Scilab plot agrees with the calculated sample coordinates
in \Cref{tab:advance}.

\subsection{Discussion}
The captured output graph moves to the left by one sample without changing any
amplitude. The sample of amplitude $3$ moves from $n=0$ to $n=-1$, and
the sample of amplitude $6$ moves from $n=1$ to $n=0$. Since the latter
is the fourth element of the vector, the output origin is four.
Delay and advance therefore use opposite changes to the origin parameter:
subtract $k$ for delay and add $k$ for advance.
